%%% chapter2.tex — Глава 2
\section{Теоретические основы сетевой стеганографии и обнаружения вредоносных вложений}

В предыдущей главе обоснована актуальность задачи обнаружения сетевой стеганографии в условиях современного ландшафта киберугроз и выделены конкретные поля заголовков TCP/IP, пригодные для встраивания скрытой информации. В настоящей главе формализуются теоретические основы как методов сокрытия данных в этих полях, так и методов их обнаружения~--- что позволяет в дальнейшем перейти к программной реализации соответствующих подсистем.

\subsection{Основные понятия и классификация методов стеганографии}

Стеганография --- это метод скрытой передачи информации путём встраивания её в другой, внешне нейтральный контейнер, с целью скрыть факт коммуникации~[2]. В отличие от криптографии, обеспечивающей конфиденциальность содержимого сообщения, стеганография скрывает само существование передаваемых данных. Сочетание обоих подходов обеспечивает максимальный уровень защиты: даже при обнаружении скрытого канала содержимое остаётся нечитаемым без ключа расшифрования.

По среде передачи стеганографические методы классифицируются следующим образом~[11]:

\begin{enumerate}[label=\arabic*)]
  \item \textbf{Текстовая стеганография} --- встраивание информации в текстовые файлы с использованием невидимых символов, пробелов, особенностей шрифтов.
  \item \textbf{Графическая стеганография} --- сокрытие данных в изображениях путём модификации младших значащих бит (LSB) пикселей, частотных компонент или цветовых каналов.
  \item \textbf{Аудио- и видеостеганография} --- кодирование данных в аудиопотоках или видеофайлах с использованием психоакустических моделей и временны\'{х} характеристик.
  \item \textbf{Сетевая стеганография} --- внедрение информации в поля и поведение сетевых протоколов при формировании и передаче пакетов TCP/IP.
\end{enumerate}

По методам встраивания выделяют три основных подхода:

\begin{enumerate}[label=\arabic*)]
  \item \textbf{С заменой (substitution)} --- внедрение информации путём замены битов или полей исходного контейнера на скрытые данные. В сетевой стеганографии этот метод применяется при модификации полей заголовков (IP~ID, TCP~ISN).
  \item \textbf{С расширением (extension)} --- добавление дополнительных полей или данных, увеличивающих объём контейнера. Например, включение дополнительных опций TCP или расширенных DNS-записей.
  \item \textbf{С генерацией (generation)} --- создание специально сформированного контейнера, изначально содержащего скрытые данные. Примером является генерация DNS-запросов с закодированными в QNAME данными.
\end{enumerate}

Формальная модель стеганографической системы включает следующие компоненты~[11, 12]:

\begin{enumerate}[label=\arabic*)]
  \item \textbf{контейнер (Cover)} --- исходный объект, в который внедряется сообщение;
  \item \textbf{стеганографический ключ} --- секретный параметр, определяющий алгоритм и место встраивания;
  \item \textbf{функция встраивания $E(C, M, K) \to S$} --- преобразование контейнера $C$, сообщения $M$ и ключа $K$ в стегоконтейнер $S$;
  \item \textbf{функция извлечения $D(S, K) \to M$} --- восстановление скрытого сообщения из стегоконтейнера;
  \item \textbf{нарушитель (Warden)} --- субъект, осуществляющий пассивный или активный анализ с целью обнаружения или уничтожения скрытой информации.
\end{enumerate}

\subsection{Сравнительный анализ сетевой и контейнерной стеганографии}

При выборе метода скрытой передачи данных по сети возникает вопрос: что целесообразнее --- встраивать информацию в поля сетевых протоколов или использовать классическую контейнерную стеганографию, передавая по сети изображения, аудио- и видеофайлы с внедрёнными данными? Оба подхода имеют свои преимущества и ограничения, и выбор между ними определяется конкретным сценарием применения.

\textbf{Контейнерная стеганография} (встраивание в изображения, аудио, видео) обладает высокой ёмкостью: метод LSB (Least Significant Bit) позволяет встроить до 12{,}5\,\% от объёма изображения (1~бит на канал на пиксель при 3~каналах), а частотные методы (DCT, DWT) обеспечивают ёмкость порядка 5--10\,\% при лучшей устойчивости к стегоанализу. Для изображения размером 1~МБ это означает возможность передачи 50--125~КБ скрытых данных в одном файле. Видеоконтейнеры ещё более ёмки: в потоке 720p~(30~fps) теоретически возможна скрытая передача данных со скоростью до нескольких Мбит/с.

Однако контейнерная стеганография при передаче по сети имеет существенные недостатки. Во-первых, необходим канал передачи контейнеров --- HTTP, FTP, электронная почта или мессенджер, что создаёт дополнительные метаданные (URL, отправитель, получатель, временны\'{е} метки), анализируемые системами DLP и SIEM. Во-вторых, промежуточные сервисы часто выполняют перекодирование мультимедиа: социальные сети сжимают изображения (JPEG с потерей качества), мессенджеры конвертируют видео, CDN-серверы оптимизируют контент. Перекодирование разрушает встроенные данные, если не использованы устойчивые (robust) алгоритмы встраивания. В-третьих, регулярная передача изображений или видеофайлов между определёнными узлами сети может привлечь внимание аналитиков безопасности, особенно если данные узлы обычно не обмениваются мультимедийным контентом.

\textbf{Сетевая стеганография} (встраивание в заголовки протоколов) обладает принципиально иными характеристиками. Её главное преимущество --- отсутствие дополнительного контейнера: скрытые данные передаются непосредственно в служебных полях пакетов, которые генерируются при любом сетевом взаимодействии. Это означает, что стеганографическая передача не создаёт новых информационных потоков и не требует дополнительных приложений или протоколов. Пакеты ICMP, DNS-запросы, TCP-соединения являются штатными элементами сетевой активности любой организации, и их наличие само по себе не вызывает подозрений.

К недостаткам сетевой стеганографии относятся низкая ёмкость каналов (от 2~до 1470~байт на пакет, что на порядки меньше, чем у контейнерной стеганографии), зависимость от сетевой инфраструктуры (NAT, межсетевые экраны, прокси-серверы могут модифицировать или отбрасывать пакеты) и более строгие требования к синхронизации между отправителем и получателем.

Таким образом, сетевая стеганография наиболее целесообразна в следующих сценариях: передача малых объёмов критически важных данных (криптографические ключи, управляющие команды, координаты C2-серверов --- типичные объёмы до нескольких килобайт); работа в средах с жёстким контролем контента (где передача мультимедийных файлов ограничена или мониторируется); необходимость полной маскировки факта коммуникации (отсутствие дополнительных файлов и соединений). Контейнерная стеганография, напротив, предпочтительна для передачи больших объёмов данных (документы, базы данных) через каналы, где мультимедийный трафик является нормой (социальные сети, мессенджеры, электронная почта).

В рамках настоящей работы выбор в пользу сетевой стеганографии обусловлен её актуальностью для задач кибербезопасности: именно сетевые скрытые каналы используются APT-группировками для управления имплантами и эксфильтрации данных в корпоративных средах с контролируемым сетевым периметром. Разработка модуля обнаружения таких каналов представляет непосредственную практическую ценность для защиты информационных систем.

\subsection{Методы сокрытия данных в заголовках сетевых протоколов}

Рассмотрим подробнее алгоритмы встраивания данных для каждого из пяти каналов, реализованных в системе \texttt{NetStego}.

\subsubsection{ICMP Payload}

Поле полезной нагрузки ICMP Echo Request (тип~8, код~0) согласно RFC~792~[5] не имеет строго определённого формата содержимого --- стандарт лишь требует, чтобы данные, отправленные в запросе, были возвращены в ответе. Это позволяет использовать до 1472~байт для встраивания скрытых данных при стандартном MTU Ethernet~[1].

Алгоритм встраивания состоит в прямом размещении данных в поле полезной нагрузки ICMP-пакета с добавлением двухбайтового префикса длины для указания объёма полезных данных. Данный канал обеспечивает наибольшую пропускную способность, однако легко обнаруживается при анализе размера и содержимого ICMP-пакетов~[11].

Основными ограничениями канала ICMP Payload являются следующие. Во-первых, многие корпоративные межсетевые экраны блокируют ICMP-трафик или ограничивают его частоту, что снижает доступность канала. Во-вторых, нестандартный размер полезной нагрузки ICMP-пакета (например, 1470~байт вместо типичных 56~или 32~байт) является явным индикатором аномалии. В-третьих, содержимое payload не шифруется на уровне протокола, и системы DPI могут выполнять содержательный анализ данных. Для повышения скрытности канала рекомендуется ограничивать размер payload значениями, характерными для конкретной ОС отправителя, а также использовать предварительное шифрование данных, что реализовано в системе \texttt{NetStego} посредством AES-256-GCM.

Граничные случаи при работе с данным каналом включают обработку пакетов с нулевой полезной нагрузкой, пакетов, превышающих MTU сети (что приводит к фрагментации на уровне IP и потенциальной потере данных при отбрасывании фрагментов), а также ситуацию, когда промежуточный маршрутизатор модифицирует содержимое ICMP payload (что возможно при использовании NAT с поддержкой ICMP ALG).

\subsubsection{IP Identification}

Поле Identification в IPv4-заголовке (16~бит) определено в RFC~791~[3] и предназначено для идентификации фрагментов датаграммы. При отсутствии фрагментации данное поле не несёт семантической нагрузки для получателя и может быть использовано для передачи 2~байт данных на пакет.

Необходимо учитывать, что операционные системы формируют значение IP~ID различными способами: инкрементально (последовательное увеличение на единицу), случайно или на основе хеш-функции. Стеганографическая модификация IP~ID может быть обнаружена путём статистического анализа распределения значений (критерий $\chi^2$)~[11].

Алгоритм встраивания в системе \texttt{NetStego} использует IP~ID для передачи 2~байт данных в каждом ICMP Echo Request пакете. Поскольку канал имеет низкую ёмкость, фреймированный поток данных разбивается на множество пакетов: для передачи 1~КБ данных требуется примерно 587~пакетов. Для обозначения конца потока используется маркерный пакет с IP~ID~=~0xFFFF, содержащий в payload двухбайтовое значение длины последнего фрагмента. Данный маркер выбран потому, что вероятность его случайного появления в нормальном трафике крайне мала (1~из 65535 для каждого пакета), а при инкрементальной генерации IP~ID значение 0xFFFF появляется лишь один раз за полный цикл.

Ограничения канала IP~ID включают невозможность передачи данных через NAT-устройства, которые могут перезаписывать поле IP~ID в процессе трансляции адресов. Кроме того, RFC~6864 допускает установку IP~ID в~0 для нефрагментируемых пакетов, и некоторые промежуточные устройства могут нормализовать данное поле.

\subsubsection{TCP Initial Sequence Number}

Начальный порядковый номер TCP-сегмента (ISN) --- 32-битное поле, устанавливаемое при инициации соединения в SYN-пакете~[4]. Согласно RFC~6528~[10], современные ОС используют криптографически стойкую рандомизацию ISN:
\begin{equation}
\text{ISN} = M + F(\text{sip}, \text{sport}, \text{dip}, \text{dport}, K),
\label{eq:isn_rfc6528}
\end{equation}
где $M$ --- счётчик с интервалом 4~микросекунды, $F$ --- псевдослучайная функция, $K$ --- секретный ключ ОС~[10]. Встраивание данных в ISN осуществляется путём полной замены значения, что требует подавления RST-пакетов ядра ОС с помощью правил межсетевого экрана.

В реализации системы \texttt{NetStego} каждый SYN-пакет содержит 4~байта скрытых данных в поле TCP~Sequence Number. Для обеспечения правильного порядка получения пакетов используются инкрементальные номера портов назначения: первый пакет направляется на порт~1024, второй --- на~1025 и т.\,д. Терминирующий пакет обозначается значением порта источника \texttt{sport~=~65535}. Данный подход позволяет получателю восстановить правильный порядок фрагментов без дополнительных метаданных.

Существенным ограничением канала TCP~ISN является генерация RST-пакетов ядром ОС в ответ на SYN-ACK от получателя. Поскольку ОС отправителя не инициировала TCP-соединение (SYN-пакет формируется библиотекой Scapy в обход TCP-стека ядра), входящий SYN-ACK воспринимается ядром как неожиданный и отвергается пакетом RST. Для предотвращения этого эффекта необходимо блокировать исходящие RST-пакеты с помощью правил межсетевого экрана (например, \texttt{iptables~-A~OUTPUT~-p~tcp~--tcp-flags~RST~RST~-j~DROP} на Linux). Данное ограничение является характерным для всех инструментов, использующих TCP~ISN для стеганографии~[20].

\subsubsection{DNS QNAME}

Протокол DNS (RFC~1035~[6]) определяет максимальную длину доменного имени в 255~байт, при этом каждая метка (label) ограничена 63~октетами. Данные кодируются в base32 и разбиваются на метки субдоменов, формируя запросы вида:
\begin{center}
\texttt{k3f7g9a1b4c...cdn-cache.net}
\end{center}

Преимуществом данного канала является его маскировка под легитимный DNS-трафик. Каждый пакет передаёт до 30~байт сырых данных при использовании двух меток по 55~символов~[6].

В реализации системы \texttt{NetStego} кодирование данных выполняется следующим образом. Входные данные разбиваются на фрагменты по 30~байт, каждый фрагмент кодируется в base32 (результат --- строка длиной до 48~символов), которая разбивается на метки по 30~символов и оформляется как субдомен конфигурируемого базового домена. Суффикс доменного имени является настраиваемым параметром (по умолчанию --- реалистичное доменное имя \texttt{cdn-cache.net}), что существенно снижает вероятность обнаружения сигнатурным методом по сравнению с использованием фиксированного суффикса. Терминирующий пакет содержит специальную метку \texttt{term} в QNAME, позволяющую получателю определить конец передачи.

Основным преимуществом канала DNS QNAME является возможность работы через корпоративные межсетевые экраны, поскольку DNS-трафик (UDP, порт~53) практически никогда не блокируется полностью. Кроме того, DNS-запросы к произвольным доменам являются штатной активностью для любой корпоративной сети. Однако данный канал имеет ряд ограничений: DNS-запросы с длинными QNAME и нечитаемыми субдоменами могут быть выявлены системами анализа DNS-трафика (DNS analytics); DNS-резолверы могут кешировать ответы, что приводит к потере повторных запросов; максимальная длина доменного имени (255~байт) ограничивает объём данных на запрос.

\subsubsection{TCP Timestamp}

Опция TCP Timestamp, описанная в RFC~7323~[7], содержит два 32-битных поля: Timestamp Value (TSval) и Timestamp Echo Reply (TSecr). Поле TSval устанавливается отправителем и используется получателем для оценки RTT и поддержки механизма PAWS~[7].

Встраивание данных осуществляется путём замены значения TSval на 4~байта скрытых данных. Для корректной работы канала необходима включённая поддержка TCP~Timestamp в операционной системе (\texttt{net.ipv4.tcp\_timestamps = 1} в Linux).

В реализации системы \texttt{NetStego} используются ACK-пакеты с установленной опцией TCP Timestamp. Каждый пакет содержит 4~байта данных в поле TSval. Для обозначения конца передачи используется маркерный пакет со значением \texttt{TSecr~=~1} (Timestamp Echo Reply), которое не встречается в нормальном TCP-трафике, поскольку TSecr является эхом TSval удалённой стороны и всегда превышает единицу в установленном соединении.

Канал TCP Timestamp обладает высокой скрытностью, поскольку опция Timestamp является стандартной частью современных TCP-реализаций и присутствует практически в каждом TCP-сегменте. Однако нарушение монотонности последовательности значений TSval является обнаруживаемой аномалией: в нормальном трафике TSval монотонно возрастает с частотой, определяемой ОС (100--1000~Гц), тогда как стеганографические значения не связаны с системным таймером и не обладают данным свойством. Для противодействия обнаружению возможно применение техники <<инкрементального кодирования>>, при которой скрытые данные добавляются к текущему значению системного таймера, однако данный подход снижает ёмкость канала и не реализован в текущей версии системы.

В таблице~\ref{tab:stego_methods_comparison} приведена сводная характеристика рассмотренных методов.

\begin{table}[htb]
\centering
\caption{Сравнение методов сокрытия данных в сетевых протоколах}
\label{tab:stego_methods_comparison}
\begin{tabular}{|l|c|c|c|c|}
\hline
Метод & Байт/пакет & Скрытность & Устойчивость & Сложность \\
 & & & к обнаружению & реализации \\
\hline
ICMP Payload   & 1470 & Низкая  & Низкая   & Низкая   \\
IP ID          & 2    & Высокая & Средняя  & Низкая   \\
TCP ISN        & 4    & Высокая & Высокая  & Высокая  \\
DNS QNAME      & 30   & Средняя & Средняя  & Средняя  \\
TCP Timestamp  & 4    & Высокая & Высокая  & Средняя  \\
\hline
\end{tabular}
\end{table}

\subsection{Методы обнаружения скрытых каналов и вредоносных вложений}

Обнаружение стеганографии в сетевых пакетах опирается на три основных направления: статистический анализ, сигнатурные методы и методы машинного обучения.

\subsubsection{Статистические методы}

Статистические методы обнаружения основаны на анализе распределений значений полей заголовков и временны\'{х} характеристик трафика.

\textbf{Энтропия Шеннона.} Вычисление энтропии поля позволяет выявить аномалии в распределении значений. Энтропия $H$ для дискретной случайной величины определяется как:
\begin{equation}
H = -\sum_{i=1}^{n} p_i \log_2 p_i,
\label{eq:shannon}
\end{equation}
где $p_i$ --- вероятность $i$-го значения в анализируемом наборе данных, $n$ --- количество уникальных значений. Высокая нормализованная энтропия (близкая к 1) свидетельствует о квазиравномерном распределении, характерном для зашифрованных или стеганографических данных~[11].

\textbf{Критерий $\chi^2$ (хи-квадрат).} Применяется для проверки гипотезы о равномерности распределения значений поля. Статистика $\chi^2$ вычисляется как:
\begin{equation}
\chi^2 = \sum_{i=1}^{k} \frac{(O_i - E_i)^2}{E_i},
\label{eq:chi_squared}
\end{equation}
где $O_i$ --- наблюдаемая частота $i$-го значения, $E_i$ --- ожидаемая частота при равномерном распределении, $k$ --- количество категорий (бинов). При $p$-значении ниже уровня значимости $\alpha$ (обычно $\alpha = 0{,}05$) гипотеза о нормальности отклоняется, что может указывать на стеганографическую модификацию.

\textbf{Критерий Колмогорова--Смирнова (KS).} Используется для сравнения распределения межпакетных интервалов (IPD) с эталонным (экспоненциальным) распределением, характерным для обычного сетевого трафика. Стеганографический трафик часто демонстрирует более регулярные интервалы~[11].

\textbf{Анализ регулярности межпакетных интервалов.} Коэффициент вариации $CV = \sigma / \mu$ межпакетных задержек позволяет оценить их регулярность. Низкий $CV$ (менее 0{,}1) свидетельствует об искусственно регулярных интервалах, характерных для автоматизированной стеганографической передачи.

\subsubsection{Сигнатурные методы}

Сигнатурные методы основаны на поиске характерных фиксированных последовательностей байт в полезной нагрузке пакетов. Примеры сигнатур:
\begin{enumerate}[label=\arabic*)]
  \item магические последовательности стеганографических протоколов (например, \texttt{NST\textbackslash x00} --- четырёхбайтовый маркер заголовка чанка системы NetStego);
  \item аномальные фиксированные значения в полях (IP~ID = 0xDEAD, TCP~TSval с нулевой энтропией);
  \item паттерн length-prefix фрейминга --- 4-байтовый префикс длины, за которым следует магическая последовательность.
\end{enumerate}

Сигнатурные методы обеспечивают высокую точность при обнаружении известных стеганографических протоколов, однако неэффективны против новых или модифицированных реализаций~[20].

\subsubsection{Методы машинного обучения}

Современный подход к обнаружению скрытых каналов использует алгоритмы машинного обучения, обученные на образцах нормального трафика~[19]. Наиболее эффективным для данной задачи является алгоритм Isolation Forest, относящийся к классу методов обнаружения аномалий (anomaly detection).

Принцип работы Isolation Forest основан на рекурсивном случайном разбиении пространства признаков. Аномалии, будучи редкими и отличными от нормальных точек, изолируются за меньшее количество разбиений. Для каждого образца $x$ вычисляется оценка аномальности~[19]:
\begin{equation}
s(x, n) = 2^{-\frac{E[h(x)]}{c(n)}},
\label{eq:isolation_forest}
\end{equation}
где $E[h(x)]$ --- математическое ожидание глубины изоляции образца $x$ по ансамблю деревьев, $c(n)$ --- средняя глубина неуспешного поиска в бинарном дереве для $n$ образцов.

Признаковое пространство для обнаружения сетевой стеганографии формируется в скользящем окне размером $w = 5$ пакетов и включает следующие признаки: значение анализируемого поля $v_i$, первая разность $\Delta v_i = v_i - v_{i-1}$, межпакетный интервал $\text{IPD}_i = t_i - t_{i-1}$, размер пакета $s_i$, скользящее среднее значения поля $\bar{v}_w = \frac{1}{w}\sum_{j=i-w+1}^{i} v_j$, скользящее стандартное отклонение значения $\sigma_{v,w}$, скользящее стандартное отклонение IPD $\sigma_{\text{IPD},w}$, энтропия payload $H_{\text{payload}}$ (бит/байт) и длина payload $l_{\text{payload}}$ (байт). Два последних признака позволяют выявлять зашифрованные данные в полезной нагрузке: зашифрованный шифротекст имеет энтропию, близкую к максимальной ($H \approx 8{,}0$~бит/байт), тогда как нормальные данные --- существенно ниже. Таким образом, для каждого пакета формируется 9-мерный вектор признаков.

Перед обучением модели все признаки нормализуются с использованием \texttt{StandardScaler}, приводящего каждый признак к нулевому среднему и единичной дисперсии. Нормализация необходима, поскольку значения полей заголовков (например, IP~ID в диапазоне 0--65535) и межпакетные интервалы (порядка 0{,}1~с) имеют существенно различные масштабы, что может негативно влиять на качество разбиения пространства в Isolation Forest.

Ключевым параметром модели является \texttt{contamination} --- ожидаемая доля аномалий в обучающей выборке. В системе \texttt{NetStego} данный параметр установлен в значение 0{,}05, что соответствует допущению о том, что не более 5\,\% обучающих данных могут содержать стеганографический трафик. Выбор данного значения обоснован экспериментальными результатами: при \texttt{contamination} $< 0{,}02$ наблюдается значительный рост ложноотрицательных срабатываний, а при \texttt{contamination} $> 0{,}1$ --- рост ложноположительных.

Количество деревьев в ансамбле (\texttt{n\_estimators}) установлено равным 100 --- значение по умолчанию библиотеки Scikit-learn~[23], обеспечивающее стабильные результаты. Размер подвыборки для построения каждого дерева (\texttt{max\_samples}) установлен в значение \texttt{auto}, что соответствует $\min(256, n)$, где $n$ --- размер обучающей выборки. Ограничение размера подвыборки является ключевым элементом алгоритма Isolation Forest, обеспечивающим как вычислительную эффективность, так и предотвращение переобучения.

Обучение модели выполняется в режиме полуавтоматического обучения (semi-supervised): модель обучается на трафике, предположительно являющемся нормальным, и затем применяется для оценки новых данных. Каждый пакет, получивший метку $-1$ (аномалия) от Isolation Forest, считается подозрительным. Итоговое решение о наличии стеганографии принимается на основе доли аномальных пакетов: если она превышает $2 \cdot \texttt{contamination} = 0{,}1$, трафик считается подозрительным с точки зрения ML-детектора.

Помимо unsupervised-подхода (Isolation Forest), в системе \texttt{NetStego} реализован supervised-классификатор на основе алгоритма Random Forest~[23]. Random Forest строит ансамбль из 100~решающих деревьев (max\_depth~=~10), каждое из которых обучается на случайной подвыборке признаков и данных. Классификатор обучается на размеченных данных: нормальный трафик (класс~0) и стеганографический трафик (класс~1). Используется тот же 9-мерный вектор признаков, что и для Isolation Forest. Преимуществом данного подхода является высокая точность классификации при наличии обучающей выборки; недостатком --- необходимость предварительной разметки данных.

\subsubsection{Анализ содержимого payload (обнаружение вредоносных вложений)}

Помимо обнаружения самого факта использования стеганографического канала, важной задачей является анализ содержимого скрытых данных для оценки уровня угрозы. Передача текстового сообщения по скрытому каналу представляет существенно меньший риск, чем передача исполняемого файла или шеллкода.

Анализ содержимого основан на сигнатурном сканировании извлечённых данных на наличие характерных последовательностей байт, свидетельствующих о присутствии определённых типов вредоносных вложений. Данный подход аналогичен методам, применяемым в антивирусных системах и системах DLP для детекции вредоносного контента в сетевом трафике.

Категории анализируемого содержимого включают: исполняемые файлы (заголовки PE и ELF), скриптовые маркеры (shebang, PowerShell, cmd.exe), индикаторы веб-атак (HTML-скрипты, вызовы eval/exec), сигнатуры шеллкода (NOP-sled), архивы (ZIP, GZIP, RAR), документы с макросами (OLE2, PDF) и обфускацию (base64-декодирование). Каждая сигнатура характеризуется уровнем уверенности, зависящим от вероятности ложноположительного срабатывания.

В таблице~\ref{tab:detection_comparison} приведено сравнение методов обнаружения.

\begin{table}[htb]
\centering
\caption{Сравнение методов обнаружения скрытых каналов}
\label{tab:detection_comparison}
\small
\begin{tabular}{|l|c|c|c|c|}
\hline
Метод & Точность & Ложные & Слож- & Обнаруж. \\
 & & срабат. & ность & каналы \\
\hline
Энтропия Шеннона     & Средняя       & Высокие      & Низкая  & ICMP, DNS \\
$\chi^2$-тест        & Средняя       & Средние      & Низкая  & IP ID, TCP TS \\
KS-тест IPD          & Средняя       & Средние      & Низкая  & Все типы \\
Isolation Forest     & Высокая       & Низкие       & Средняя & Все типы \\
Random Forest        & Высокая       & Низкие       & Средняя & Все (superv.) \\
Сигнатуры            & Оч. высокая   & Оч. низкие   & Низкая  & Известные \\
Анализ содержимого   & Оч. высокая   & Низкие       & Низкая  & Payload \\
\hline
\end{tabular}
\end{table}

\subsection{Обзор существующих программных решений}

Существующие инструменты обнаружения вторжений обладают базовыми возможностями для детекции классических атак, однако специализированных механизмов для выявления стеганографии в сетевых протоколах недостаточно~[12].

Среди систем обнаружения вторжений следует выделить:

\begin{enumerate}[label=\arabic*)]
  \item \textbf{Zeek (Bro)} --- фреймворк анализа сетевого трафика с поддержкой пользовательских скриптов. Позволяет реализовать детекцию стеганографии через механизм событий, однако требует значительных усилий по разработке и не содержит готовых модулей для этой цели.
  \item \textbf{Snort и Suricata} --- системы обнаружения вторжений на основе сигнатур и правил. Поддерживают написание пользовательских правил для анализа полей пакетов, но не обеспечивают встроенной статистической или ML-детекции.
\end{enumerate}

Среди стеганографических инструментов известны:

\begin{enumerate}[label=\arabic*)]
  \item \textbf{Hydan} --- исследовательский инструмент для сокрытия данных в TCP~ISN~[20]. Ограничен одним каналом, не поддерживает шифрование.
  \item \textbf{Nstego} --- пакетный стеганографический инструмент для ICMP и IP-полей. Не обеспечивает фрагментацию и контроль целостности.
  \item \textbf{StegTunnel} --- реализация стеганографического VPN-канала. Ориентирован на туннелирование, а не на скрытую передачу файлов.
\end{enumerate}

Общим недостатком перечисленных аналогов является отсутствие комплексного подхода, объединяющего криптографическую защиту, фрагментацию с контролем целостности, многоканальную передачу и многоуровневое обнаружение в единой системе.

\subsection{Подходы, используемые при анализе сетевого трафика}

Для анализа сетевого трафика и реализации как стеганографических каналов, так и модуля обнаружения используются два основных подхода.

\textbf{PCAP и live capture.} Формат PCAP (Packet Capture) позволяет анализировать ранее сохранённый трафик без необходимости привилегий администратора. Live capture с использованием библиотеки libpcap предоставляет возможность перехвата и анализа трафика в реальном времени, но требует повышенных привилегий (CAP\_NET\_RAW на Linux, административные права на Windows)~[1].

\textbf{BPF-фильтры.} Berkeley Packet Filter (BPF) обеспечивает эффективную фильтрацию пакетов на уровне ядра ОС. Примеры фильтров для различных каналов:
\begin{enumerate}[label=\arabic*)]
  \item \texttt{icmp and icmp[0] == 8} --- ICMP Echo Request;
  \item \texttt{udp port 53} --- DNS-трафик;
  \item \texttt{tcp[tcpflags] \& tcp-syn != 0} --- TCP SYN-пакеты.
\end{enumerate}

\textbf{Библиотеки для работы с пакетами.} Для реализации системы \texttt{NetStego} выбрана библиотека Scapy~[21], предоставляющая возможности создания, модификации, отправки и анализа сетевых пакетов произвольной структуры. В отличие от dpkt (ориентированного на скорость парсинга) и PyShark (обёртки над tshark), Scapy обеспечивает полный контроль над формированием полей заголовков, что критически важно для реализации стеганографических каналов.
